% =============================================================================
% Template LaTeX IPG  —  https://github.com/VagnerBomJesus/TEMPLATE-IPG-LATEX
% Autor base, criador e mantenedor do projeto:
%     Vagner Bom Jesus  (@VagnerBomJesus)
%
% Proprietario e autor original deste template. Ao reutilizar, distribuir ou
% editar, MANTER este credito ao autor base (nao remover). Quem usa o template
% e o autor do SEU relatorio; o criador base do modelo e Vagner Bom Jesus.
% Licenca: GPL-3.0 (ver LICENSE).
% =============================================================================

% =============================================================================
% Preâmbulo — Pacotes e Configurações
% IMPORTANTE: requer XeLaTeX ou LuaLaTeX (NÃO pdflatex)
% =============================================================================

% --- Língua ---
\usepackage[portuguese]{babel}

% --- Fontes (requer XeLaTeX/LuaLaTeX + fontspec) ---
\usepackage{fontspec}
\setmainfont{Times New Roman}
\setsansfont{Calibri}[
    BoldFont       = Calibri Bold,
    ItalicFont     = Calibri Italic,
    BoldItalicFont = Calibri Bold Italic
]
\setmonofont{Courier New}

% --- Geometria da página ---
\usepackage[a4paper,top=2.5cm,bottom=2.5cm,inner=3cm,outer=2.5cm]{geometry}

% --- Espaçamento entre linhas ---
\usepackage{setspace}
\onehalfspacing

% --- Headings (titlesec) ---
\usepackage{titlesec}
\titleformat{\chapter}[block]{\sffamily\bfseries\fontsize{18}{22}\selectfont}{\thechapter.}{0.5em}{}
\titlespacing*{\chapter}{0pt}{0pt}{18pt}
\titleformat{\section}{\sffamily\bfseries\fontsize{16}{20}\selectfont}{\thesection}{0.5em}{}
\titlespacing*{\section}{0pt}{18pt}{6pt}
\titleformat{\subsection}{\sffamily\bfseries\fontsize{14}{18}\selectfont}{\thesubsection}{0.5em}{}
\titlespacing*{\subsection}{0pt}{14pt}{4pt}
\titleformat{\subsubsection}{\sffamily\fontsize{12}{15}\selectfont}{\thesubsubsection}{0.5em}{}
\titlespacing*{\subsubsection}{0pt}{12pt}{4pt}
\titleformat{\paragraph}[hang]{\sffamily\fontsize{12}{15}\selectfont}{\theparagraph}{0.5em}{}
\titlespacing*{\paragraph}{0pt}{10pt}{4pt}

% --- Cabeçalhos e rodapés — Calibri 10pt ---
\usepackage{fancyhdr}
\pagestyle{fancy}
\fancyhf{}
\fancyhead[LE]{\sffamily\fontsize{10}{12}\selectfont\leftmark}
\fancyhead[RO]{\sffamily\fontsize{10}{12}\selectfont\rightmark}
\fancyfoot[C]{\sffamily\fontsize{10}{12}\selectfont\thepage}
\renewcommand{\headrulewidth}{0.4pt}
\renewcommand{\footrulewidth}{0pt}
\fancypagestyle{plain}{\fancyhf{}\fancyfoot[C]{\sffamily\fontsize{10}{12}\selectfont\thepage}\renewcommand{\headrulewidth}{0pt}}

% --- Gráficos e imagens ---
\usepackage{graphicx}
\graphicspath{{images/}}

% --- Cores ---
\usepackage{xcolor}
\definecolor{azulIPG}{HTML}{00467F}

% --- Tabelas ---
\usepackage{booktabs}
\usepackage{array}
\usepackage{longtable}

% --- Matemática ---
\usepackage{amsmath}
\usepackage{amssymb}

% --- Código fonte (listings) — tudo a preto ---
\usepackage{listings}
\lstset{
    basicstyle=\ttfamily\small,
    keywordstyle=\bfseries,
    commentstyle=\itshape,
    stringstyle=\color{black},
    numbers=left,
    numberstyle=\tiny,
    numbersep=8pt,
    breaklines=true,
    frame=single,
    framexleftmargin=15pt,
    captionpos=b,
    tabsize=4,
    showstringspaces=false
}

% --- Hiperligações + metadados do PDF ---
\usepackage{hyperref}
\hypersetup{
    colorlinks=true,
    linkcolor=black,
    citecolor=black,
    urlcolor=black,
    bookmarksnumbered=true,
    pdfauthor={\nomeAutor},
    pdftitle={\tituloRelatorio},
    pdfcreator={Template LaTeX IPG criado por Vagner Bom Jesus (@VagnerBomJesus) -- github.com/VagnerBomJesus/TEMPLATE-IPG-LATEX},
    pdfproducer={XeLaTeX -- Template base de Vagner Bom Jesus (@VagnerBomJesus)}
}

% --- Aspas (recomendado pelo biblatex; carregar antes) ---
\usepackage{csquotes}

% =============================================================================
% ESTILO DAS REFERÊNCIAS — deixar UM bloco ativo e comentar os outros com %.
% Depois de trocar: apagar auxiliares e recompilar (xelatex->biber->xelatex x2).
% =============================================================================

% ---- ATIVO: IEEE (numérico [1], por ordem de citação) ----
\usepackage[
    backend=biber,
    style=ieee,
    sorting=none,
    natbib=true,
    maxbibnames=99
]{biblatex}

% ---- OPÇÃO: APA 7.ª ed. (requer o pacote biblatex-apa) ----
% \usepackage[backend=biber,style=apa,sorting=nyt,natbib=true,maxbibnames=99]{biblatex}

% ---- OPÇÃO: Autor-data / Harvard (NP 405) ----
% \usepackage[backend=biber,style=authoryear,sorting=nyt,natbib=true,maxbibnames=99]{biblatex}

% ---- OPÇÃO: Numérico simples ----
% \usepackage[backend=biber,style=numeric-comp,sorting=nyt,natbib=true,maxbibnames=99]{biblatex}

\addbibresource{backmatter/referencias.bib}

% --- Legendas — Calibri 10pt, separador traço médio (–) ---
\usepackage{caption}
\DeclareCaptionLabelSeparator{travessao}{ \textendash\ }
\captionsetup{font={sf,small},labelfont={sf,bf,small},format=hang,labelsep=travessao}

% --- Subcapturas ---
\usepackage{subcaption}

% --- Indentação do primeiro parágrafo ---
\usepackage{indentfirst}
\setlength{\parindent}{1.25cm}

% --- Enumerações e listas ---
\usepackage{enumitem}

% --- PDFs nos anexos ---
\usepackage{pdfpages}

% --- Numeração por capítulo ---
\usepackage{chngcntr}
\counterwithin{figure}{chapter}
\counterwithin{table}{chapter}
\counterwithin{equation}{chapter}
% lstlisting só existe no início do documento -> adiar com \AtBeginDocument
\AtBeginDocument{\counterwithin{lstlisting}{chapter}}

% --- Profundidade do índice ---
\setcounter{tocdepth}{3}
\setcounter{secnumdepth}{4}

% --- Glossário / Acrónimos ---
\usepackage[acronym,nonumberlist,toc]{glossaries}
\makeglossaries
